\section*{Chapter \ref{ch:g10:what-prices-do} --- What Prices Do}
% figdata: figdata/grade-10/09-what-prices-do/out/ (fuel_rates, hair_rates, hicp_summary); Brenland numbers invented.

\begin{solution}{exo:g10:what-prices-do:1}
Informs: coffee drinkers, roasters and cafés learn that coffee is scarcer without knowing about the drought. Reason to act: drinkers switch some cups to tea,
growers in other countries harvest more carefully and plant more. Rations: the coffee left goes to the buyers willing to pay most for it.
\end{solution}

\begin{solution}{exo:g10:what-prices-do:2}
Between 28.2\,\% (Italy) and 40.3\,\% (France). The hairdressing lines never fell below zero (the lowest value is 0.0, in Spain in early 2014): a negative value would
mean that a haircut cost less than a year earlier.
\end{solution}

\begin{solution}{exo:g10:what-prices-do:3}
$100 - 2p = 3p - 50$ gives $p^* = 30$, $Q^* = 40$. At 25: demanded $100 - 50 = 50$, supplied $75 - 50 = 25$, a shortage of 25.
\end{solution}

\begin{solution}{exo:g10:what-prices-do:4}
Demanded 40, supplied 20: a shortage of 20 (the equilibrium is 30 units at 30).
\end{solution}

\begin{solution}{exo:g10:what-prices-do:5}
25 units are traded instead of 40. At 25 units the demand price is $(100 - 25)/2 = 37.5$ and the supply price 25. The \omterm{def:g10:what-prices-do:dwl}{deadweight loss} is $\tfrac12 \times (40 - 25) \times
(37.5 - 25) = \tfrac12 \times 15 \times 12.5 = 93.75$ dalens.
\end{solution}

\begin{solution}{exo:g10:what-prices-do:6}
20 units are traded instead of 30; at 20 units the demand price is 40 and the supply price 20. \omterm{def:g10:what-prices-do:dwl}{Deadweight loss} $\tfrac12 \times 10 \times 20 = 100$. On the 20 units sold,
buyers pay 10 less than at equilibrium: $20 \times 10 = 200$ dalens move from sellers to buyers.
\end{solution}

\begin{solution}{exo:g10:what-prices-do:7}
Supplied $40 - 24 = 16$, demanded $120 - 80 = 40$: a shortage of 24. At 16 units the demand price is $30 - 4 = 26$ and the supply price 20: \omterm{def:g10:what-prices-do:dwl}{deadweight loss} $\tfrac12 \times (24 -
16) \times 6 = 24$. Moved to buyers: $16 \times (24 - 20) = 64$ dalens.
\end{solution}

\begin{solution}{exo:g10:what-prices-do:8}
Demanded $100 - 72 = 28$, supplied $108 - 50 = 58$: an \omterm{def:g10:market-equilibrium:excess-demand}{excess supply} of 30. At 28 units the supply price is $(28 + 50)/3 = 26$ and the demand price 36: \omterm{def:g10:what-prices-do:dwl}{deadweight loss}
$\tfrac12 \times (40 - 28) \times (36 - 26) = 60$ dalens.
\end{solution}

\begin{solution}{exo:g10:what-prices-do:9}
Buy at 200, carry for 20, sell at 230: a gain of 10 a tonne. Traders do so until the price has risen in the first town and fallen in the second enough for the gap to equal
20. With a cost of 40, the gap of 30 is too small: nobody carries wheat and the gap stays.
\end{solution}

\begin{solution}{exo:g10:what-prices-do:10}
Buying at 10, storing for 5 and selling at 18 gains 3 a sack: traders buy and store. Their buying raises the harvest price and the extra rice they sell later lowers the lean
season price. Storage stops growing when the gap has fallen to the storage cost, 5 dalens.
\end{solution}

\begin{solution}{exo:g10:what-prices-do:11}
Up to 40 dalens. The 20 units go to buyers whose values are 40 or more; the next buyer, who gets nothing, values a unit at just under 40, and would pay nearly that much to get
one. So a holder can sell a unit on for about 40, the demand price at 20 units.
\end{solution}

\begin{solution}{exo:g10:what-prices-do:12}
A gap between the prices of the same good in two places that stays larger than the cost of moving it, with no rule stopping trade. Before concluding, check that the goods are
really the same (quality, freshness), that all costs are counted (taxes, insurance, losses on the way, the time it takes), that the prices are for the same day, and that traders
can learn both prices.
\end{solution}

\begin{solution}{exo:g10:what-prices-do:13}
The equilibrium price is 24, below the ceiling of 26: the ceiling allows every trade that happens anyway, and nothing changes. With $Q = 138 - 4p$: $138 - 4p = 2p - 24$ gives $6p =
162$, $p = 27$, $Q = 30$. The ceiling now binds: sellers offer $52 - 24 = 28$, buyers want $138 - 104 = 34$, a shortage of 6.
\end{solution}

\begin{solution}{exo:g10:what-prices-do:14}
Each of the 12 buyers loses 9 dalens of time: $12 \times 9 = 108$ dalens destroyed. \omterm{def:g10:surplus-and-efficiency:consumer-surplus}{Consumer surplus} falls from 126 to 18, and \omterm{def:g10:surplus-and-efficiency:total-surplus}{total surplus} to $18 + 36 = 54$, a quarter of the
216 at equilibrium.
\end{solution}

\begin{solution}{exo:g10:what-prices-do:15}
At 36, buyers take 28 and sellers offer 58: the government buys 30 units for $30 \times 36 = 1080$ dalens. Sellers' surplus rises from $\tfrac12 \times 40 \times (30 - \tfrac{50}{3})
\approx 266.7$ to $\tfrac12 \times 58 \times (36 - \tfrac{50}{3}) \approx 560.7$, a gain of 294; buyers' surplus falls from 400 to $\tfrac12 \times 28 \times 14 = 196$. The taxpayers
pay 1080 for a gain to sellers of 294.
\end{solution}

\begin{solution}{pb:g10:what-prices-do:1}
Rents in dalens a month.

\textbf{1.} $10000 - 5p = 4000 + p$ gives $6p = 6000$: $p^* = 1000$, $Q^* = 5000$ flats.

\textbf{2.} Offered $4000 + 800 = 4800$, wanted $10000 - 4000 = 6000$: a shortage of 1200.

\textbf{3.} The last tenant would pay $(10000 - 4800)/5 = 1040$. The last landlord rents out at 800 ($4800 = 4000 + p$).

\textbf{4.} $\tfrac12 \times (5000 - 4800) \times (1040 - 800) = \tfrac12 \times 200 \times 240 = \text{\dal{24000}}$ a month.

\textbf{5.} $4800 \times (1000 - 800) = \text{\dal{960000}}$ a month.

\textbf{6.} $1040 - 800 = 240$ dalens a month.

\textbf{7.} $800 + 240 = 1040$, more than the 1000 of the free market. Landlords would get the surplus, and more: the side payments bring them $4800 \times 240 = \text{\dal{1152000}}$ a
month above the legal rent.

\textbf{8.} At most 240 dalens a month per flat, $4800 \times 240 = \text{\dal{1152000}}$ a month.

\textbf{9.} With more tenants than flats, a landlord who neglects repairs still finds a tenant: lower quality is another way of raising the true price.

\textbf{10.} $Q = 5000 + 10(p - 1000) = 10p - 5000$.

\textbf{11.} $8000 - 5000 = 3000$ flats; the shortage is $6000 - 3000 = 3000$.

\textbf{12.} $(10000 - 3000)/5 = 1400$ dalens.

\textbf{13.} $\tfrac12 \times (5000 - 3000) \times (1400 - 800) = \tfrac12 \times 2000 \times 600 = \text{\dal{600000}}$ a month.

\textbf{14.} $600000/24000 = 25$ times.

\textbf{15.} $4800 - 3000 = 1800$ flats fewer than in the first year, and \num{2000} fewer than at equilibrium (40\,\%). In San Francisco controlled landlords cut their rental housing by
15\,\%; the direction is the same.

\textbf{16.} Short run: $1 \times 1000/5000 = 0.2$. Long run: $10 \times 1000/5000 = 2$.

\textbf{17.} First year $4800 \times 200 = \text{\dal{960000}}$; after five years $3000 \times 200 = \text{\dal{600000}}$ a month.

\textbf{18.} The 3000 tenants who hold flats gain 200 a month each. Landlords lose. Of the 6000 households who want a flat at 800, 3000 find none; \num{2000} households who would
rent at the free-market rent of 1000 can no longer do so.

\textbf{19.} Long-run supply at 950: $9500 - 5000 = 4500$. Demand price at 4500: $(10000 - 4500)/5 = 1100$. \omterm{def:g10:what-prices-do:dwl}{Deadweight loss} $\tfrac12 \times 500 \times (1100 - 950) = \text{\dal{37500}}$ a month.

\textbf{20.} With the rent capped at 800 dalens, the \omterm{def:g10:what-prices-do:dwl}{deadweight loss} is \dal{24000} a month in the first year and \dal{600000} after five years; the black-market rent rises from 1040 to
1400 dalens.
\end{solution}
